\chapter{Estimating a Normal Population Mean}
\label{chap:normal-mean}


This worked example shows how to connect a statistical model, an estimator,
an uncertainty statement, and an interpretation. The observations are
illustrative; they do not represent a scientific study. General treatments of
estimation and confidence intervals are given by \citet{wasserman2004all}.


\section{Model and estimand}
Assume $X_1,\ldots,X_n$ are independent normal random variables with mean
$\mu$ and variance $\sigma^2$, where $n\geq 2$ and $\sigma^2>0$. The target is
$\mu$, the population mean. For observed values $x_1,\ldots,x_n$, the density
of one observation is
\begin{equation}
  \label{eq:normal-density}
  f(x;\mu,\sigma^2)
  = \frac{1}{\sqrt{2\pi\sigma^2}}
    \exp\left\{-\frac{(x-\mu)^2}{2\sigma^2}\right\},
  \qquad x\in\RR.
\end{equation}
Equation~\eqref{eq:normal-density} specifies the sampling model, not a
probability distribution on the unknown parameter $\mu$.


\section{Maximum likelihood estimation}
The log likelihood is
\begin{equation}
  \label{eq:normal-loglikelihood}
  \ell(\mu,\sigma^2)
  = -\frac{n}{2}\log(2\pi\sigma^2)
    -\frac{1}{2\sigma^2}\sum_{i=1}^{n}(x_i-\mu)^2.
\end{equation}
Holding $\sigma^2$ fixed in Equation~\eqref{eq:normal-loglikelihood},
differentiation with respect to $\mu$ gives
\begin{align*}
  \frac{\partial\ell}{\partial\mu}
    &= \frac{1}{\sigma^2}\sum_{i=1}^{n}(x_i-\mu), \\
  \frac{\partial^2\ell}{\partial\mu^2}
    &= -\frac{n}{\sigma^2}<0.
\end{align*}
The score is zero at $\widehat{\mu}=\bar{x}$. The negative second derivative
shows that this is the unique maximizer in $\mu$ for each fixed variance.
If the observed values are not all equal, joint maximization also gives
$\widehat{\sigma}^2_{\mathrm{ML}}
= n^{-1}\sum_{i=1}^{n}(x_i-\bar{x})^2$. This variance estimate differs from
the unbiased sample variance used below. For identical observations, the
joint likelihood has no finite maximizer with $\sigma^2>0$.


\section{Confidence intervals}
\begin{prop}[Normal-mean pivot]
  \label{prop:normal-pivot}
  Under the normal sampling model, define
  \[
    S^2 = \frac{1}{n-1}\sum_{i=1}^{n}(X_i-\overline{X})^2.
  \]
  Then
  \begin{equation}
    \label{eq:normal-pivot}
    T = \frac{\overline{X}-\mu}{S/\sqrt{n}}\sim t_{n-1}.
  \end{equation}
\end{prop}
\begin{proof}
  The standardized mean is standard normal, $(n-1)S^2/\sigma^2$ has a
  chi-squared distribution with $n-1$ degrees of freedom, and these two
  quantities are independent. Their ratio, with the chi-squared quantity
  divided by its degrees of freedom and square-rooted, has the stated
  Student distribution.
\end{proof}
Let $t_{\nu,p}$ denote the $p$ quantile of a Student distribution with $\nu$
degrees of freedom. Proposition~\ref{prop:normal-pivot} and
Equation~\eqref{eq:normal-pivot} give a $100(1-\alpha)\%$ confidence interval
for $\mu$, for $0<\alpha<1$:
\begin{equation}
  \label{eq:mean-confidence-interval}
  \left[\bar{x}-t_{n-1,1-\alpha/2}\frac{s}{\sqrt{n}},\quad
        \bar{x}+t_{n-1,1-\alpha/2}\frac{s}{\sqrt{n}}\right].
\end{equation}
Here $s$ is the observed sample standard deviation. Appendix~\ref{app:variance}
explains the $n-1$ denominator in the sample variance.


For the illustrative values $1,2,3$, we have $n=3$, $\bar{x}=2$, and $s=1$.
Using $t_{2,0.975}\approx 4.303$ in
Equation~\eqref{eq:mean-confidence-interval} gives a 95\% confidence interval
of approximately $[-0.484,4.484]$. The wide interval reflects the uncertainty
associated with three observations. These values demonstrate a calculation;
they do not provide evidence for the normal sampling assumption.


\begin{remark}[Coverage]
  The interval procedure covers the fixed population mean in 95\% of repeated
  samples under the stated model. A realized frequentist interval does not
  assign a 95\% probability to the fixed parameter. Exact coverage relies on
  normality and independence; a small sample alone does not justify a normal
  approximation.
\end{remark}


\section{Planning a simulation study}
A simulation can examine interval coverage when the sampling model changes.
Before generating data, specify the aim, data-generating mechanisms, estimand,
methods, and performance measures \citep{morris2019using}. For example,
compare normal and skewed populations with the same mean and variance, vary
sample size, and evaluate the interval in
Equation~\eqref{eq:mean-confidence-interval}.


Report the repetition count, Monte Carlo uncertainty, and failed fits.
Retain reproducible code. This template reports no simulation results.
